\section{Continuous Assessment 1 (CA-1) --- Set B}
\label{sec:ca1_set_b}

\LPUPaperHeader{Continuous Assessment (CA-1)}{B}{60 Minutes}{30 Marks}{CO1 (Matrix Algebra), CO2 (Homogeneous Differential Equations)}

\begin{center}
\sffamily\textbf{\color{AlertRed}Note: Attempt all six questions. Each question carries 5 marks.}
\end{center}

\questionitem{1}{
    Find the canonical normal form and rank of the matrix:
    \[
    A = \begin{bmatrix} 2 & 3 & 4 & 5 \\ 3 & 4 & 5 & 6 \\ 4 & 5 & 6 & 7 \end{bmatrix}
    \]
}{5}{CO1}{L3 Apply}

\questionitem{2}{
    Test the linear independence of the functions $y_1 = \sin x$, $y_2 = \cos x$ using the Wronskian.
}{5}{CO2}{L2 Understand}

\questionitem{3}{
    Find the eigenvalues of $A^3$ and $A^{-1}$ if $A = \begin{bmatrix} 1 & 4 \\ 2 & 3 \end{bmatrix}$.
}{5}{CO1}{L3 Apply}

\questionitem{4}{
    Solve $(D^4 - 16)y = 0$.
}{5}{CO2}{L3 Apply}

\questionitem{5}{
    Solve the initial value problem: $(D^2 + 4D + 4)y = 0$, $y(0) = 1$, $y'(0) = -1$.
}{5}{CO2}{L3 Apply}

\questionitem{6}{
    Apply the Cayley-Hamilton theorem to find $A^{-1}$ for $A = \begin{bmatrix} 1 & 3 \\ 2 & 4 \end{bmatrix}$.
}{5}{CO1}{L3 Apply}

\newpage
\subsection*{Solutions and Marking Rubrics --- CA-1 Set B}

\begin{solutionbox}[Solution to Question 1]
Row operations: $R_2 \to R_2 - R_1$, $R_3 \to R_3 - R_2$ gives rows of 1s. A subsequent subtraction reduces row 3 to zero. Normal form is $\begin{bmatrix} I_2 & 0 \\ 0 & 0 \end{bmatrix}$. Thus, $\mathbf{\operatorname{rank}(A) = 2}$.
\end{solutionbox}

\begin{solutionbox}[Solution to Question 2]
The Wronskian is:
\[
W(\sin x, \cos x) = \begin{vmatrix} \sin x & \cos x \\ \cos x & -\sin x \end{vmatrix} = -\sin^2 x - \cos^2 x = -1 \neq 0
\]
Since $W \neq 0$ for all $x \in \mathbb{R}$, the functions are \textbf{Linearly Independent}.
\end{solutionbox}

\begin{solutionbox}[Solution to Question 3]
$\det(A - \lambda I) = \lambda^2 - 4\lambda - 5 = 0 \implies \lambda = 5, -1$.
Eigenvalues of $A^3$ are $\lambda^3$: $5^3 = 125$ and $(-1)^3 = -1$.
Eigenvalues of $A^{-1}$ are $\frac{1}{\lambda}$: $\frac{1}{5}$ and $-1$.
\end{solutionbox}

\begin{solutionbox}[Solution to Question 4]
$m^4 - 16 = (m^2 - 4)(m^2 + 4) = (m - 2)(m + 2)(m^2 + 4) = 0 \implies m = \pm 2, \pm 2i$.
Solution: $\mathbf{y(x) = c_1 e^{2x} + c_2 e^{-2x} + c_3 \cos(2x) + c_4 \sin(2x)}$.
\end{solutionbox}

\begin{solutionbox}[Solution to Question 5]
$(m + 2)^2 = 0 \implies m = -2, -2$.
$y(x) = (c_1 + c_2 x)e^{-2x}$. $y(0) = c_1 = 1$.
$y'(0) = -2c_1 + c_2 = -1 \implies c_2 = 1$.
Unique solution: $\mathbf{y(x) = (1 + x)e^{-2x}}$.
\end{solutionbox}

\begin{solutionbox}[Solution to Question 6]
$\det(A - \lambda I) = \lambda^2 - 5\lambda - 2 = 0 \implies A^2 - 5A - 2I = 0 \implies A^{-1} = \frac{1}{2}(A - 5I) = \frac{1}{2}\begin{bmatrix} -4 & 3 \\ 2 & -1 \end{bmatrix}$.
\end{solutionbox}

\newpage
